\akapit Prace laboratoryjne zostały podzielone na kilka niezależnych etapów badawczych, realizowanych przy użyciu przygotowanego środowiska symulacyjnego[cite: 8]:
\begin{enumerate}
    \item \textbf{Lot swobodny:} Wykonano loty sterowane manualnie w celu wyznaczenia warunków równowagi modelu (prędkości, kąta natarcia oraz wychylenia steru wysokości) dla zadanych warunków[cite: 8].
    \item \textbf{Badanie dynamiki serwomechanizmów:} Przeprowadzono automatyczne loty według dwóch zaprogramowanych trajektorii dla wariantu referencyjnego oraz z uwzględnieniem trzech różnych prędkości pracy serwomechanizmów[cite: 8, 9].
    \item \textbf{Badanie luzów w układzie sterowania:} Wykonano serię lotów automatycznych dla trzech różnych poziomów luzów kinematycznych w torze sterowania[cite: 9].
    \item \textbf{Badanie opóźnień transportowych:} Zrealizowano loty automatyczne przy włączonej symulacji opóźnień dla trzech różnych wartości czasu opóźnienia sygnału sterującego[cite: 10].
    \item \textbf{Badanie szumów pomiarowych:} Przeprowadzono symulacje lotu automatycznego z aktywnym modelem szumów czujników, badając wpływ trzech różnych wartości mocy szumu pomiarowego na stabilność i dokładność śledzenia trasy[cite: 10, 11].
\end{enumerate}
\akapit W każdym z powyższych wariantów badawczych rejestrowano kluczowe zmienne stanu, sygnały zadane oraz sygnały symulowane w celu późniejszej analizy numerycznej i wyznaczenia wskaźników jakościowych[cite: 8, 12, 13].
